\documentclass[12pt]{article}
\usepackage[spanish]{babel}
\usepackage[letterpaper,top=2cm,bottom=2cm,left=3cm,right=3cm,marginparwidth=1.75cm]{geometry}
\usepackage{tabularx}
\usepackage{fancyvrb}
\usepackage{graphicx}
\usepackage[document]{ragged2e}

\renewcommand{\baselinestretch}{1.25}

\setlength{\parindent}{0pt}

\DefineVerbatimEnvironment{Json}{Verbatim}{fontsize=\small, frame=single, framerule=0.5pt, framesep=5pt}

\begin{document}

\begin{titlepage}

    \centering

    \vspace*{1 cm}
    
    \Huge
    \textbf{Proyecto: Cafetería - Documentación Técnica}
    
    \vspace{0.5 cm}
    
    \Large
    Programación de Servicios y Procesos
    
    \vspace{5.5 cm}
    
    \textbf{Adrián Condines Celada}
    
    \vspace{0.8 cm}
        
    \normalsize
    Aula Estudio\\

    \vspace{0.8 cm}

    2º Ciclo Superior - Desarrollo de Aplicaciones Multiplataforma\\

    \vspace{0.8 cm}

    Curso 2025 - 2026

\end{titlepage}

\tableofcontents
\newpage

\section{Introducción y Objetivos}

\justify
El objetivo de esta práctica es desarrollar una aplicación en Java que simule el funcionamiento de una cafetería en la que los clientes llegarán periódicamente a la cafetería a realizar su pedido y serán atendidos por los camareros disponibles. 

La aplicación debe gestionar la llegada de clientes, la atención por parte de los camareros y la salida de los clientes una vez que han sido atendidos.

\section{Requisitos Funcionales}

\justify
La aplicación debe cumplir con los siguientes requisitos funcionales: \newline
RF-1: El sistema debe permitir la creación de objetos del tipo \texttt{Cliente} con los atributos \texttt{nombre} y \texttt{tiempoEspera}.\newline
RF-2: Cada cliente debe ejecutarse en un hilo independiente que simule su llegada a la cafetería.\newline
RF-3: Cada cliente debe intentar pedir un café y esperar un tiempo determinado \texttt{(tiempoEspera)}.\newline
RF-4: Si el café no está preparado antes de que se cumpla su \texttt{tiempoEspera}, el cliente se irá y se reflejará con un mensaje.\newline
RF-5: Si el cliente recibe su café a tiempo, se reflejará con un mensaje que incluya que el cliente ha sido servido y cuanto tiempo ha tomado en ser atendido.\newline
RF-6: El sistema debe permitir la creación de objetos del tipo \texttt{Camarero}.\newline
RF-7: Cada camarero debe ejecutarse en un hilo independiente que simule la atención a los clientes.\newline
RF-8: Los camareros deben preparar los pedidos de los clientes en base al orden de llegada.\newline
RF-9: Si un cliente se va antes de que su café esté listo, el camarero debe cancelar la preparación del café y pasar al siguiente cliente en la cola.\newline
RF-10: El programa debe de seguir ejecutándose hasta que todos los clientes hayan sido atendidos o se hayan ido.\newline
RF-11: Se deben mostrar mensajes en la consola que indiquen el estado de los clientes y camareros en cada momento. 

\section{Requisitos No Funcionales}

\justify
La aplicación debe cumplir con los siguientes requisitos no funcionales: \newline
RNF-1: El programa debe de manejar múltiples hilos de ejecución para simular la concurrencia entre clientes y camareros.\newline
RNF-2: Se debe garantizar que cada cliente sea atendido una sola vez.\newline
RNF-3: Los mensajes en la consola deben ser claros y proporcionar información relevante sobre el estado de la cafetería.\newline
RNF-4: El código debe estar bien estructurado y documentado para facilitar su comprensión y mantenimiento.\newline

\section{Casos de Uso}

\justify
A continuación se describen los casos de uso principales de la aplicación:

\section{Historias de Usuario}

\justify
A continuación se describen las historias de usuario principales de la aplicación:

\section{Estructura del Proyecto}

\justify
La estructura del proyecto está organizada de la siguiente manera:



\section{Clases Principales}

\subsection{Clase Cliente}

\subsubsection{Descripción}
La clase \texttt{Cliente} representa a un cliente que llega a la caferería para pedir un café.

\subsubsection{Atributos}
Los atributos de la clase \texttt{Cliente} son los siguientes:

\begin{itemize}

    \item \texttt{String} \textbf{nombre}: Nombre del cliente.
    \item \texttt{String} \textbf{tiempoEspera}: Tiempo máximo que el cliente está dispuesto a esperar por su café.
    \item \texttt{boolean} \textbf{atendido}: Indica si el cliente ha sido atendido o no.
    \item \texttt{long} \textbf{inicioEspera}: Marca de tiempo que indica cuándo comenzó a esperar el cliente.
    \item \texttt{Queue<Cliente>} \textbf{colaClientes}: Cola compartida de clientes esperando ser atendidos.

\end{itemize}

\subsubsection{Métodos}

\subsection{Clase Camarero}

\subsubsection{Descripción}
La clase \texttt{Camarero} reprsenta a un camarero que atiende a los clientes en la cafetería.

\subsubsection{Atributos}

\begin{itemize}

    \item \texttt{String} \textbf{nombre}: Nombre del camarero.
    \item \texttt{int} \textbf{cafesServidos}: Contador de cafés servidos por el camarero.
    \item \texttt{Queue<Cliente>} \textbf{colaClientes}: Cola compartida de clientes esperando ser atendidos.

\end{itemize}

\subsubsection{Métodos}

\end{document}